\documentclass[10pt]{article}
\usepackage[margin=0.82in]{geometry}
\usepackage{amsmath,amssymb,amsthm,mathtools,booktabs,array,microtype,hyperref}
\hypersetup{hidelinks,pdftitle={An exact obstruction to the quadratic half-line completion}}
\newtheorem{theorem}{Theorem}
\newtheorem{lemma}[theorem]{Lemma}
\newtheorem{proposition}[theorem]{Proposition}
\newcommand{\R}{\mathbb R}
\newcommand{\V}{\mathcal V_9}
\newcommand{\Cq}{\mathcal C_{33}}
\newcommand{\Cr}{\mathcal C_{\rm ram}}
\newcommand{\ip}[2]{\langle #1,#2\rangle}
\newcommand{\tr}{\operatorname{tr}}
\newcommand{\diag}{\operatorname{diag}}
\newcommand{\sgn}{\operatorname{sgn}}
\newcommand{\co}{\operatorname{coeff}}
\title{An exact obstruction to the quadratic half-line completion}
\author{}
\date{}
\begin{document}
\maketitle
\vspace{-2.5em}
\begin{abstract}
The proposed order-33 cone is a proper subcone of the full nonnegative
canonical valid-slack cone. We give an all-positive rational source and an
explicit rational separating functional, certified by four positive definite
dual matrices. A fixed enlargement, obtained by allowing a selector quartic
in the square root of the Bob variable, contains the excluded source. Its
order-49 affine pencil and exact feasible matrices are supplied. The two certificates isolate the role of the additional square-root terms.
\end{abstract}

\section{The two cones and the counterexample}
Use the source order
\[
c=(d,\lambda,n,m,C,A,D,\alpha,\beta)\in\R_+^9,
\qquad D_u=1+u^2,\quad D_v=1+v^2.
\]
The full canonical cone $\V$ consists of the sources for which the canonical numerator
\begin{align}
F_c={}&4\lambda D_u(W+vX)^2+4nD_uD_vX^2
 +4m(vW-X+uvY-uZ)^2\notag\\
&+4CD_v(uW-Y)^2+4AD_uD_vY^2
 +(4D+2d)(uW+uvX-Y-vZ)^2\notag\\
&+(4\alpha+2d)D_u(vY-Z)^2
 +(4\beta+2d)D_v(X+uZ)^2-2dD_uD_vZ^2                       \label{eq:F}
\end{align}
is nonnegative for every real state vector and every pair of measurement
angles, including their projective endpoints. Complex states are included
by separating real and imaginary parts of this real symmetric quadratic form.

The state reduction gives four quadratic matrices and two completions:
\begin{align}
P&=\begin{pmatrix}4mv^2&-4mv\\-4mv&4m+(4\beta+2d)D_v\end{pmatrix},&
N&=\begin{pmatrix}4CD_v+4D+2d&(4D+2d)v\\(4D+2d)v&(4D+2d)v^2\end{pmatrix},\notag\\
E&=\begin{pmatrix}4AD_v+(4\alpha+2d)v^2&-(4\alpha+2d)v\\
-(4\alpha+2d)v&4\alpha-2dv^2\end{pmatrix},&
K&=4\lambda\begin{pmatrix}1&v\\v&v^2\end{pmatrix}
       +4nD_v\begin{pmatrix}0&0\\0&1\end{pmatrix},            \label{eq:PNEK}\\
G_J&=\begin{pmatrix}P+K&J\\J^T&N+K\end{pmatrix},&
H_J&=\begin{pmatrix}P+E&K-E-J\\K-E-J^T&N+E\end{pmatrix}.  \label{eq:GH}
\end{align}
Let $\Cq$ be the cone obtained by requiring a quadratic selector
$J(v)=J_0+vJ_1+v^2J_2$ on $v\geq0$. Equivalently, writing
$Q(v)=Q_0+vQ_1+v^2Q_2$ for $Q=G_J,H_J$, its four PSD blocks are
\begin{equation}
S_Q=\begin{pmatrix}Q_0&B_Q\\B_Q^T&Q_2\end{pmatrix}\succeq0,
\qquad T_Q=Q_1-B_Q-B_Q^T\succeq0.                         \label{eq:primal}
\end{equation}
The nine source-sign blocks give total order $8+8+4+4+9=33$ and 44
auxiliaries. The two-LMI identity gives $\Cq\subseteq\V$.
Reflection by $T=\diag(1,-1)$ supplies the negative half-line; homogeneous
continuity supplies the projective endpoints. No restriction on state signs
is made.

\begin{theorem}\label{thm:main}
The all-positive rational source
\begin{equation}
\boxed{\bar c=
 \left(1,2^{-34},2^{-14},2^{-24},\frac1{576},161,
                     \frac12,76,\frac1{250}\right)}         \label{eq:cbar}
\end{equation}
belongs to $\V$ but not to $\Cq$. Moreover, the fixed quartic-in-$\sqrt v$
completion cone $\Cr$ defined in Section~\ref{sec:enlargement} satisfies
\[
\boxed{\Cq\subsetneqq\Cr\subseteq\V,\qquad
       \bar c\in\Cr\setminus\Cq.}
\]
It has an evaluated rational affine pencil of order 49 with nine source
coordinates and 148 auxiliaries.
\end{theorem}

\section{The exact separating inequality}
Here is the integer linear functional that separates \eqref{eq:cbar}:
\begin{align}
\ell(c)={}&-23803618d+906393494300170\lambda+2744379904n\notag\\
 &+1378590360912m+5935589516C+20425A\notag\\
 &+2600028D+91248\alpha+278104928\beta.                    \label{eq:ell}
\end{align}
Its exact value is
\begin{equation}
\boxed{\ell(\bar c)=-\frac{5417798011632274183}{9663676416000}
                    <-560635.}                            \label{eq:negative}
\end{equation}
The dual matrices below prove $\ell\geq0$ on $\Cq$.

\begin{lemma}[Dual trace identity]\label{lem:dual}
For $\sigma=G,H$, let $U_\sigma,R_\sigma,W_\sigma$ be symmetric $4\times4$
matrices satisfying
\begin{equation}
M_\sigma=\begin{pmatrix}U_\sigma&R_\sigma\\R_\sigma&W_\sigma\end{pmatrix}
\succeq0,\qquad R_\sigma\succeq0.                         \label{eq:dualPSD}
\end{equation}
For each of the three pairs $(U_G,U_H)$, $(R_G,R_H)$, $(W_G,W_H)$, require
their upper-right $2\times2$ blocks to agree. These are twelve scalar
equations. Put $M_{\sigma,0}=U_\sigma$, $M_{\sigma,1}=R_\sigma$,
$M_{\sigma,2}=W_\sigma$. Then
\begin{equation}
f(c)=\sum_{k=0}^2\left[\ip{M_{G,k}}{G_k(c,0)}
                         +\ip{M_{H,k}}{H_k(c,0)}\right]    \label{eq:f}
\end{equation}
is nonnegative on $\Cq$. Here $G_k(c,0),H_k(c,0)$ mean the coefficients
with $J=0$, and $\ip AB=\tr(AB)$.
\end{lemma}
\begin{proof}
For each side, expand \eqref{eq:primal} to obtain
\[
\ip{M_\sigma}{S_\sigma}+\ip{R_\sigma}{T_\sigma}
=\ip{U_\sigma}{Q_0}+\ip{R_\sigma}{Q_1}+\ip{W_\sigma}{Q_2}.
\]
The unrestricted $B_\sigma$ cancels because $R_\sigma$ is symmetric.
The coefficients $J_k$ enter $G_k,H_k$ through opposite off-diagonal blocks,
so the twelve overlap equations cancel every entry of $J_0,J_1,J_2$.
Consequently
\begin{equation}
f(c)=\ip{M_G}{S_G}+\ip{R_G}{T_G}
     +\ip{M_H}{S_H}+\ip{R_H}{T_H}\geq0.                  \label{eq:pairing}
\end{equation}

\end{proof}

Appendix~\ref{app:dual} displays six rational $4\times4$ matrices satisfying
\eqref{eq:dualPSD} and the twelve overlap equations, with
\begin{equation}
                         f(c)=2^{-42}\ell(c).              \label{eq:scale}
\end{equation}
Their four PSD blocks are in fact positive definite. The source-sign dual
variables are all zero. Equations \eqref{eq:negative}--\eqref{eq:scale}
therefore prove $\bar c\notin\Cq$ by finite rational arithmetic.

This also excludes a polynomial half-line selector of any degree. Indeed,
every $2\times2$ principal minor of $G_J$ gives
\[
|J_{ij}(v)|^2\leq(P+K)_{ii}(v)(N+K)_{jj}(v)=O(v^4)
\qquad(v\longrightarrow+\infty).
\]
Every polynomial entry must therefore have degree at most two, already
excluded by the certificate above. Every selector for this source therefore depends nonpolynomially on $v$.

\section{Validity of the source}
Validity follows from the rational contact source
\begin{equation}
\begin{aligned}
(d,\lambda,n)&=(1,0,0),\\
(m,C,A)&=\left(\frac{1166523000}{31661820415017653},
\frac{322803000}{186199784281},\frac{218288946200}{1359122513}\right),\\
(D,\alpha,\beta)&=\left(\frac{90943434293}{222507771414},
\frac{248802430236}{3300726103},\frac{205476973}{51646655494}\right),
\end{aligned}                                               \label{eq:cstar}
\end{equation}
denoted $c_*$. Its exact contact is
$(u,v,\psi)=(1/13,16,(140,-1/32,1/50,1))$.
Let $\Pi_*(u,v)$ be the
$4\times4$ state matrix of \eqref{eq:F}, and set
$x=169u^2$, $y=v^2/256$. The retained determinant identity can be written
\begin{equation}
\det\Pi_*=16\kappa D_u^2D_v^2u^2v^2\,\mathcal A(x,y),
\quad
\kappa=\frac{365997345792000000}
 {806534883104527170621493981142180739537045858457}>0,        \label{eq:det}
\end{equation}
where
\begin{align*}
\mathcal A(x,y)={}&\rho_1(x+x^{-1}-2)+\rho_2(y+y^{-1}-2)\\
&+\rho_3(y/x+x+y^{-1}-3)+\rho_4(xy+x^{-1}+y^{-1}-3),
\end{align*}
and the positive integer coefficients are
\[
\begin{array}{c|r}
\rho_1&1267917833244593541025180323231\\
\rho_2&762274793296209851400745596627083\\
\rho_3&2737791187307806620394799996019\\
\rho_4&79165812050999420348852524500
\end{array}
\]
Each bracket is nonnegative by the two- or three-term AM--GM inequality.
To pass from determinant positivity to all-state positivity, the $Z=0$
principal block $B_*$ satisfies
\begin{equation}
B_*\succeq T_*^T\diag(4m,4CD_v,4AD_uD_v)T_*\succ0,
\qquad
T_* =\begin{pmatrix}v&-1&uv\\u&0&-1\\0&0&1\end{pmatrix},
\quad\det T_*=u                                           \label{eq:transverse}
\end{equation}
when $u\ne0$. All omitted terms of \eqref{eq:F} are nonnegative on $Z=0$.
Thus the scalar Schur complement and \eqref{eq:det} give
$\Pi_*\succeq0$ for $uv\ne0$. Polynomial continuity recovers the axes;
bihomogenization in the two angle pairs recovers every projective endpoint.


Direct rational comparison gives $\bar c-c_*\geq0$, with the $d$ coordinate
unchanged. Every other coordinate in \eqref{eq:F} multiplies a nonnegative
event form. Therefore
\[
                       F_{\bar c}=F_{c_*}+F_{\bar c-c_*}\geq0
\]
for all states and tests. This proves $\bar c\in\V$. The unrounded contact
source itself is also separated:
\begin{equation}
\ell(c_*)=-\frac{303722700569565123211387668040}
                    {267603838100324717114473}<0.           \label{eq:oldnegative}
\end{equation}

\section{One fixed enlargement that repairs the obstruction}
\label{sec:enlargement}
Set $v=16s^2$ on the positive half-line. Allow a selector
$J(s)=\sum_{k=0}^4J_ks^k$. This adds the selector powers $\sqrt v$ and
$v^{3/2}$, with auxiliary coefficients independent of the testing variable.
Put
\[
Z_2(s)=\begin{pmatrix}I_4\\sI_4\\s^2I_4\end{pmatrix},
\qquad Z_1(s)=\begin{pmatrix}I_4\\sI_4\end{pmatrix}.
\]
For $Q=G_J,H_J$, impose the identity
\begin{equation}
Q(16s^2)=Z_2(s)^TA_QZ_2(s)+sZ_1(s)^TB_QZ_1(s),
\qquad A_Q\succeq0,\quad B_Q\succeq0.                    \label{eq:ramGram}
\end{equation}
Here the notation on the left means substitution in the source blocks,
with $J=J(s)$. The block orders are $12,8,12,8$ and nine scalar source
blocks. This defines $\Cr$ and immediately proves $\Cr\subseteq\V$
using the two-LMI identity. Matrix half-line factorization also
shows that \eqref{eq:ramGram} captures exactly the quartic-in-$s$ selector
condition, including singular matrices. Substituting $v=16s^2$ in any
quadratic half-line certificate gives one of these Grams, proving
$\Cq\subseteq\Cr$.

\paragraph{Eliminating the coefficient equations.}
For a prescribed quartic $Q(s)=\sum_{k=0}^4Q_ks^k$, write
$A=(A_{ij})_{i,j=0}^2$ and $B=(B_{ij})_{i,j=0}^1$ in $4\times4$ blocks.
Choose $B_{00},B_{11}$ symmetric, $B_{01},A_{02}$ arbitrary, and
$K_{01},K_{12}$ skew-symmetric. Set
\begin{align}
A_{00}&=Q_0,& A_{22}&=Q_4,\notag\\
A_{01}&=\tfrac12(Q_1-B_{00})+K_{01},&
A_{12}&=\tfrac12(Q_3-B_{11})+K_{12},\notag\\
A_{11}&=Q_2-A_{02}-A_{02}^T-B_{01}-B_{01}^T.               \label{eq:param}
\end{align}
Transpose the lower blocks. These formulas parametrize every coefficient
solution of \eqref{eq:ramGram}. They use $10+10+16+16+6+6=64$ variables
per side, in addition to the 20 selector coefficients. Thus the final
pencil has 148 auxiliaries and no coefficient equations. The exported
49-dimensional pencil has 411 nonzero upper-triangular rational coefficient
entries, with every source and auxiliary coefficient evaluated.

\paragraph{Exact repair of both sources.}
The supplied data give rational face matrices $V_i$ and positive definite
rational matrices $X_i$ of orders $7,3,11,7$. Their congruences give the
four Gram matrices in \eqref{eq:ramGram} at $c_*$. Every unpivoted LDL
pivot of $X_i$ is greater than $2^{-18}$; all 80 coefficient conditions
after selector elimination are checked exactly. The full rational Grams for
$\bar c$ are also exported, together with their 148 affine-pencil
auxiliaries. Their PSD proof is the same four positive congruences plus
nonnegative rational event outer products.

The extension from $c_*$ to $\bar c$ is elementary. For $\delta=\bar c-c_*$,
one has $d_\delta=0$, and $P_\delta,N_\delta,E_\delta,K_\delta$ are sums of
polynomial event outer products. Taking $J_\delta=K_\delta$ gives
\begin{align*}
G_\delta&=\diag(P_\delta,N_\delta)
    +\begin{pmatrix}K_\delta&K_\delta\\K_\delta&K_\delta\end{pmatrix},\\
H_\delta&=\diag(P_\delta,N_\delta)
    +\begin{pmatrix}E_\delta&-E_\delta\\-E_\delta&E_\delta\end{pmatrix}.
\end{align*}
These additions lie in the same quartic-in-$s$ family. This also proves that
the entire $d=0$ source face is included, with arbitrary zero event weights.
For example, the new selector at $c_*$ has the nonzero odd coefficient
\[
[s]J_{21}(s)=
\frac{444900885879842604203466822371109}
     {390926447097959500000000000000000}.
\]
The explicit certificate uses the additional square-root terms. The separation proves that some
function outside the quadratic span is necessary.

\paragraph{The fixed scalar testing family.}
To express the Gram certificate directly in the physical state
$\psi=(W,X,Y,Z)^T$, set
\[
\eta=L(u)\psi,
\quad
L(u)=\begin{pmatrix}u&0&u^2&0\\0&u&0&u^2\\u&0&-1&0\\0&u&0&-1\end{pmatrix},
\quad U=u^2,\quad D_U=\diag(1,1,U,U).
\]
For $v\ne0$ put $s=\sqrt{|v|/16}$ and
$\eta_\sigma=\diag(1,\sigma,1,\sigma)\eta$, where $\sigma=\sgn(v)$.
At $v=0$ choose $\sigma=1$. The four groups of fixed functions are
\begin{equation}
\boxed{
Z_2(s)D_U\eta_\sigma,\quad
\sqrt{s}\,Z_1(s)D_U\eta_\sigma,\quad
uZ_2(s)\eta_\sigma,\quad
u\sqrt{s}\,Z_1(s)\eta_\sigma.}                            \label{eq:features}
\end{equation}
They have sizes $12,8,12,8$ and their block Gram form equals $u^2D_uF_c$.
The Bob factors are exactly
$\{1,s,s^2,\sqrt{s},s^{3/2}\}$, so the added Gram factors are fixed quarter
powers of $|v|$. Source and testing variables remain independent.
Dividing every function by $D_u^{3/2}\sqrt{D_v}$ represents
\begin{equation}
\frac{u^2}{D_u}\,\frac{F_c}{D_uD_v}.                       \label{eq:normalized}
\end{equation}
This division involves only testing variables. Here is a literal compact
algebraic cover, including every endpoint. Take
\[
a_0^2+a_1^2=1,\quad a_0\geq0,\qquad
\rho^2+\tau^2=1,\quad\rho,\tau\geq0,\qquad\epsilon\in\{-1,1\},
\]
and put
\[
\kappa=\sqrt{\tau^8+256\rho^8},\quad
b=\frac{(\tau^4,16\epsilon\rho^4)}{\kappa},\quad
g=\frac{(\tau^4,\rho^2\tau^2,\rho^4)}{\kappa},\quad
r=\frac{(\rho\tau^3,\rho^3\tau)}{\kappa}.
\]
These are source-independent algebraic functions with
$\kappa^2\geq\tau^8+\rho^8\geq1/8$.
In the interiors, $u=a_1/a_0$, $v=16\epsilon(\rho/\tau)^4$.
Define the homogeneous cubic vectors
\begin{align*}
P_G&=\left(\begin{smallmatrix}
a_0^2a_1W+a_0a_1^2Y\\
\epsilon(a_0^2a_1X+a_0a_1^2Z)\\
a_1^3W-a_0a_1^2Y\\
\epsilon(a_1^3X-a_0a_1^2Z)
\end{smallmatrix}\right),&
P_H&=\left(\begin{smallmatrix}
a_0a_1^2W+a_1^3Y\\
\epsilon(a_0a_1^2X+a_1^3Z)\\
a_0a_1^2W-a_0^2a_1Y\\
\epsilon(a_0a_1^2X-a_0^2a_1Z)
\end{smallmatrix}\right).
\end{align*}
The forty functions on this closed cover are
\[
                 g\otimes P_G,\quad r\otimes P_G,
                 \quad g\otimes P_H,\quad r\otimes P_H.
\]
Their Gram sum is $a_1^2f_c(a,b;\psi)$, where $f_c$ is the bihomogeneous
normalized physical slack. This follows from \eqref{eq:normalized} on the
interior and is an algebraic identity everywhere. Each function has absolute
value less than $3$ when $\|\psi\|=1$. At Bob's endpoints,
$(g,r)=((1,0,0),0)$ or $((0,0,1/16),0)$; the Alice endpoints are literal
polynomial substitutions in $P_G,P_H$. The multiplier $a_1^2$ is positive
on a dense testing set. Consequently every state sign and every endpoint
is retained in the validity converse.

Exchanging Alice and Bob gives an
\href{../ramified-obstruction/Q222_Ramified_Quartic_Obstruction.pdf}{all-positive valid source outside $\Cr$}.
Thus $\Cr\subsetneq\V$. Increasing the polynomial degree in $s$ cannot
repair this exclusion: the PSD principal-minor bounds force every
polynomial selector to have degree at most four.

\appendix
\section{Six rational matrices proving the separation}\label{app:dual}
Put $q=2^{30}$. The following are the hatted moment matrices. The lower
triangles are displayed as well so that there is no storage convention
ambiguity.
\begingroup
\small
\allowdisplaybreaks
\setlength{\arraycolsep}{3pt}
\begin{align*}
q\widehat U_G&=\begin{pmatrix}
3376576&-4595&8890485/2^{20}&-141059/2^{14}\\
-4595&82702&374101/2^7&-2966\\
8890485/2^{20}&374101/2^7&48657&-1435456\\
-141059/2^{14}&-2966&-1435456&159900862
\end{pmatrix},\\[3pt]
q\widehat R_G&=\begin{pmatrix}
2570&373&100461333/2^{23}&-796959/2^{16}\\
373&5619&132105/2^5&-4192\\
100461333/2^{23}&132105/2^5&4654&-50824\\
-796959/2^{16}&-4192&-50824&1442242
\end{pmatrix},\\[3pt]
q\widehat W_G&=\begin{pmatrix}
20&26&156439171/2^{23}&-18103/2^{10}\\
26&822&92585/2^7&-4442\\
156439171/2^{23}&92585/2^7&697&-4948\\
-18103/2^{10}&-4442&-4948&53164
\end{pmatrix},\\[3pt]
q\widehat U_H&=\begin{pmatrix}
142545776&374885&142247760&-564236\\
374885&1002&374101&-1483\\
142247760&374101&141950400&-563057\\
-564236&-1483&-563057&2251
\end{pmatrix},\\[3pt]
q\widehat R_H&=\begin{pmatrix}
201343592&529527&200922666&-796959\\
529527&1401&528420&-2096\\
200922666&528420&200502635&-795292\\
-796959&-2096&-795292&3163
\end{pmatrix},\\[3pt]
q\widehat W_H&=\begin{pmatrix}
313472891&94149&312878342&-1158592\\
94149&16684&92585&-2221\\
312878342&92585&312285067&-1156239\\
-1158592&-2221&-1156239&4509
\end{pmatrix}.
\end{align*}
\endgroup

Set
\[
\Delta_G=\diag(128,1/32,1/2,1/32),\qquad
\Delta_H=\diag(1/512,1/16,1/512,1/32).
\]
The moments used in Lemma~\ref{lem:dual} are
\[
U_\sigma=\Delta_\sigma\widehat U_\sigma\Delta_\sigma,
\quad R_\sigma=16\Delta_\sigma\widehat R_\sigma\Delta_\sigma,
\quad W_\sigma=256\Delta_\sigma\widehat W_\sigma\Delta_\sigma.
\]
The twelve upper-right-block equalities follow entry by entry. For positive
definiteness, apply the rational recurrence
\[
p_j=A_{jj}-\sum_{k<j}L_{jk}^2p_k,\qquad
L_{ij}=\frac{A_{ij}-\sum_{k<j}L_{ik}L_{jk}p_k}{p_j},\quad L_{jj}=1.
\]
For $\widehat M_\sigma=\left(\begin{smallmatrix}
\widehat U_\sigma&\widehat R_\sigma\\
\widehat R_\sigma&\widehat W_\sigma\end{smallmatrix}\right)$,
the exact floors of $2^{30}p_j$ are
\[
\begin{array}{c|l}
\widehat M_G&(3376576,82695,48553,117468710,16,14,14,31)\\
\widehat R_G&(2570,5564,1593,10646)\\
\widehat M_H&(142545776,16,32,17,29078313,12,24,11)\\
\widehat R_H&(201343592,8,15,8).
\end{array}
\]
The first scaled pivot in every row equals the listed integer; each later
pivot lies strictly between its listed integer and that integer plus one.
In particular, every pivot is greater than $2^{-27}$.
Finally, direct coefficient pairing with \eqref{eq:PNEK} gives
\eqref{eq:scale}. These six small matrices, the stated recurrence, and one
rational substitution prove the separation.

\section{Closedness and strict feasibility of the ramified cone}
\paragraph{An integral bound and closedness.}
For indices starting at zero, put
\[
 (H_3)_{ij}=\frac1{i+j+1}\quad(0\le i,j\le2),\qquad
 (H'_2)_{ij}=\frac1{i+j+2}\quad(0\le i,j\le1),
\]
and $W=\operatorname{diag}(H_3\otimes I_4,H'_2\otimes I_4,
H_3\otimes I_4,H'_2\otimes I_4)$. This is positive definite: its factors
are the moment matrices of $1,s,s^2$ and of $1,s$ with weights $1$ and
$s$ on $[0,1]$. Their inverse traces are $381$ and $54$, respectively,
so $W\succeq I_{40}/381$. The completion terms cancel in the trace, giving
the exact coefficient identity
\begin{equation}
 \left\langle W,\operatorname{diag}(A_G,B_G,A_H,B_H)\right\rangle
 =\int_0^1\operatorname{Tr}(G_J(16s^2)+H_J(16s^2))\,ds
 =8\int_0^1(1+256s^4)\,ds\sum_{k=0}^8c_k
 =\frac{2088}{5}\sum_{k=0}^8c_k.                    \label{eq:ramified-moment}
\end{equation}
Thus $W$ annihilates all $148$ auxiliary coefficient matrices.
If all four Gram blocks are positive semidefinite, bounded source
coordinates bound their total trace and hence every entry. The auxiliary
variables are recovered by fixed linear operations: the $J_k$ are the
upper-right $2\times2$ blocks of the coefficients of $G_J(16s^2)$;
$B_{00},B_{11},B_{01},A_{02}$ are Gram subblocks; and $K_{01},K_{12}$
are the skew parts of $A_{01},A_{12}$. Bounded sources therefore admit
bounded feasible auxiliary witnesses. Taking a subsequence proves that
the projected source cone is closed. For a strictly positive weight
on the whole pencil, use $\widetilde W=\operatorname{diag}(I_9,W)$;
it still annihilates the auxiliaries, and its nine source contractions
are all $2093/5$.

\paragraph{A strict primal point.}
The rational order-$33$ certificate in
\href{https://github.com/anubhav70747/q222-finite-sdp/blob/main/extracted-lifts/canonical-completion/full-polar/slater_proof.txt}{the accompanying strict-feasibility proof} has
$c^0=(1/100,1,1,1,1,1,1,1,1)$ and $J(v)=K(v)$.
For either of its completion polynomials $Q(v)=Q_0+vQ_1+v^2Q_2$,
write its established certificate as
\[
 S=\begin{pmatrix}Q_0&B\\B^T&Q_2\end{pmatrix},\qquad
 T=Q_1-B-B^T=I_4/50.
\]
The two old bounds are $S_G\succeq19I_8/25$ and
$S_H\succeq179I_8/300$. Consequently
\[
 A^0=\begin{pmatrix}
 Q_0&0&16B\\0&16T&0\\16B^T&0&256Q_2
 \end{pmatrix}\succeq\frac8{25}I_{12}.
\]
Set $B^{\rm new}=I_8/100$, and replace the $01$ and $12$ blocks of
$A^0$, and their transposes, by $-I_4/200$. These changes cancel
the additional coefficients $sI_4/100+s^3I_4/100$ from
$sZ_1^TB^{\rm new}Z_1$, preserving the exact polynomial $Q(16s^2)$.
The symmetric perturbation has maximum absolute row sum $1/100$.
Hence $A^{\rm new}\succeq31I_{12}/100$ for both $Q=G,H$.
The full order-$49$ pencil is therefore at least $I_{49}/100$.

In the actual $148$-coordinate parametrization, this point has
\[
 (J_0,J_1,J_2,J_3,J_4)=
 \left(4I_2,0,
       \begin{pmatrix}0&64\\64&0\end{pmatrix},0,
       \begin{pmatrix}0&0\\0&2048\end{pmatrix}\right),
\]
and, on each side,
$B_{00}=B_{11}=I_4/100$, $B_{01}=0$, $A_{02}=16B$,
$K_{01}=K_{12}=0$. The literal auxiliary vector and the four Gram
blocks are supplied in the accompanying rational data.

\end{document}
